\chapter{Introducción}

El presente proyecto, titulado \textit{“Análisis, Diseño y Construcción de un Sistema Operativo desde Cero”}, se desarrolla en el marco del curso de Sistemas Operativos de la Escuela Profesional de Ingeniería Informática y de Sistemas de la Universidad Nacional de San Antonio Abad del Cusco. El curso tiene como propósito fortalecer la comprensión de los principios, componentes y decisiones de diseño que sustentan a un sistema operativo; en consecuencia, este proyecto se plantea como un ejercicio académico orientado a aprender mediante investigación, análisis comparativo y fundamentación técnica.

El documento se organiza en \textbf{tres partes}. La primera parte desarrolla la base conceptual y teórica: conceptos fundamentales, clasificación de sistemas operativos, arquitecturas y componentes esenciales. La segunda parte amplía el análisis hacia proyectos reales de sistemas operativos desarrollados desde cero, priorizando iniciativas que no se apoyen directamente en una base Unix preexistente, con el fin de evaluar enfoques, herramientas, complejidad y valor formativo. La tercera parte consolida los hallazgos y formaliza la \textbf{selección de un sistema de referencia}, el cual fue \textbf{ejecutado en un entorno de emulación/virtualización} y \textbf{documentado} como caso de estudio, relacionando sus componentes y comportamiento con los contenidos abordados en el curso.

En conjunto, este documento unificado integra la evidencia recolectada, organiza criterios de comparación y expone la justificación técnica y pedagógica de la propuesta seleccionada. Si bien el proyecto se titula con énfasis en la construcción desde cero, en la fase documentada en este informe no se implementa un sistema operativo propio; el foco se centra en el análisis y en la documentación aplicada sobre un sistema existente seleccionado.

\section{Objetivos de la investigación}

\subsection{Objetivo General}
Analizar y documentar fundamentos, tipologías, arquitecturas y componentes esenciales de los sistemas operativos mediante un estudio comparativo de propuestas desarrolladas desde cero, con el fin de sustentar la selección de un sistema de referencia y establecer un marco técnico que guíe la documentación aplicada y la planificación de un trabajo de diseño posterior.

\subsection{Objetivos Específicos}
\begin{itemize}
    \item Describir conceptos fundamentales, funciones y evolución de los sistemas operativos, así como su importancia en la gestión de recursos y la ejecución de programas.
    \item Examinar clasificaciones y arquitecturas relevantes (monolítica, microkernel e híbrida), identificando sus ventajas, limitaciones e implicancias en rendimiento, seguridad y mantenibilidad.
    \item Investigar y comparar un conjunto representativo de sistemas operativos desarrollados desde cero, identificando propósito, arquitectura, lenguajes de implementación, componentes, herramientas de construcción, documentación disponible y complejidad estimada.
    \item Sistematizar los hallazgos mediante tablas, esquemas y figuras que permitan un análisis comparativo claro y verificable.
    \item Definir criterios técnicos y pedagógicos de selección, y justificar la elección de SerenityOS como sistema de referencia para el proyecto.
    \item Ejecutar el sistema seleccionado en un entorno controlado (emulación/virtualización) y documentar, con enfoque académico, aspectos vinculados a los contenidos del curso.
    \item Establecer lineamientos iniciales del entorno de trabajo y una planificación de alto nivel para etapas posteriores de diseño.
\end{itemize}

\section{Alcance del análisis}
El alcance de la investigación se delimita a los siguientes aspectos:
\begin{itemize}
    \item \textbf{Ámbito conceptual:} comprende la revisión de fundamentos teóricos y técnicos, incluyendo clasificación, arquitectura y componentes principales de los sistemas operativos.
    \item \textbf{Ámbito comparativo:} incluye el análisis de sistemas operativos desarrollados desde cero (con énfasis en propuestas no derivadas directamente de Unix), considerando propósito, arquitectura, lenguajes, herramientas, documentación y viabilidad académica.
    \item \textbf{Ámbito de selección:} abarca la definición de criterios y la justificación del sistema de referencia para orientar la documentación aplicada y etapas posteriores de diseño.
    \item \textbf{Ámbito documental:} se limita a fuentes públicas y verificables (documentación oficial, repositorios, artículos y referencias técnicas), evitando afirmaciones no sustentadas.
\end{itemize}

\section{Limitaciones}
\begin{itemize}
    \item La investigación se apoya principalmente en análisis documental y comparativo; en consecuencia, la disponibilidad, actualidad y precisión de la documentación consultada condicionan la profundidad de algunos apartados.
    \item Los sistemas operativos analizados presentan alcances heterogéneos (didácticos, experimentales o de propósito general), lo que limita comparaciones directas en ciertos criterios y exige contextualizar resultados.
    \item Se considera únicamente información pública y verificable; no se incluyen detalles que dependan de fuentes privadas, internas o no accesibles.
    \item Este documento se limita a la selección de un sistema de referencia, su ejecución en entornos controlados y la documentación de hallazgos relevantes para el curso.
\end{itemize}

\section{Metodología}
La metodología aplicada es de carácter documental y comparativa, y comprende las siguientes etapas:
\begin{itemize}
    \item \textbf{An\'alisis de información:} recopilación de referencias académicas y técnicas, documentación oficial, repositorios y materiales de apoyo asociados a cada sistema operativo analizado.
    \item \textbf{Definición de criterios:} establecimiento de criterios técnicos y pedagógicos (arquitectura, lenguajes, herramientas, documentación, complejidad y viabilidad de aprendizaje) para orientar el análisis.
    \item \textbf{Sistematización:} organización de la información en tablas, esquemas y figuras que permitan comparar de manera consistente los criterios definidos.
    \item \textbf{Análisis comparativo:} evaluación de similitudes y diferencias entre propuestas, identificando patrones, fortalezas y limitaciones relevantes para el contexto académico.
    \item \textbf{Selección y justificación:} elección de un sistema de referencia y argumentación de la decisión con base en la evidencia recolectada y los criterios definidos.
    \item \textbf{Ejecución y documentación aplicada:} ejecución del sistema seleccionado en un entorno de emulación/virtualización y registro de observaciones técnicas vinculadas a los contenidos del curso.
\end{itemize}